% Compile from this directory with: xelatex -no-shell-escape notes-template.tex
% Compile twice to resolve references. Uses local, licensed EB Garamond fonts.
\documentclass[11pt]{article}
\usepackage{cadearticle}
% Edit these document fields. Leave optional values empty to omit their lines.
% short-title is optional; the running header uses title when it is blank.
\notesetup{
  title={Faith and Promise},
  subtitle={Reading Romans 4 in light of the Augsburg Confession},
  author={Cade Ferreira},
  date={12 October 2026},
  source={Personal study notes · Romans 4},
  short-title={Faith and Promise}
}
\begin{document}
\noteopening

\section{The Promise}
Paul locates the promise in the faithfulness of God. Faith receives what has been promised; it does not supply a second foundation for the promise. These sample paragraphs demonstrate the note template rather than transcribe a published article.

This distinction shapes the way assurance is understood. The reader is directed to the promise itself, rather than to the strength of the act of believing. The same layout can hold reading notes, an article draft, or a sustained reflection on a passage.

\begin{bquote}
That is why it depends on faith, in order that the promise may rest on grace.\footnote{Romans 4:16, ESV (excerpt); compare Augsburg Confession IV.}
\end{bquote}

\subsection{Faith and Assurance}
The source and the response belong together, but they do not perform the same task. This paragraph illustrates the first-line indent, justified text, and close paragraph spacing of the supplied Cade Article style.

\section{Questions for Further Study}
How does Paul connect the promise to Abraham with the life of the church? What should a reader attend to when considering the relationship between faith, grace, and assurance?

\subsection{Sources to Revisit}
Return to Romans 4 and Galatians 3. Record bibliographic details in footnotes so that the note remains useful when it is read separately from the application.\footnote{This is sample template content. Replace the title, author, source line, headings, paragraphs, quotations, and references with your own material.}
\end{document}
